\section{Backend Architecture}

The backend of ByteVault is a high-performance HTTP API built with \textbf{Bun} and \textbf{Express}. It utilizes TypeScript for type safety and \texttt{zod} for strict runtime schema validation.

\subsection{Architectural Pattern}
The application recently transitioned from a tightly-coupled Router-Controller model to a \textbf{Layered Service Architecture}. This isolates business logic, database queries, and HTTP transport layers.

\begin{itemize}
    \item \textbf{Routes}: Define HTTP endpoints and apply middleware.
    \item \textbf{Middleware}: Handle cross-cutting concerns like Auth (\texttt{requireEmployeeAuth}) and Idempotency.
    \item \textbf{Services}: Encapsulate reusable business logic (e.g., \texttt{AccountService}, \texttt{AuditService}, \texttt{LedgerService}).
\end{itemize}

\textbf{Why it is there:} A layered architecture prevents catastrophic failures. For instance, if the Audit log database insert fails, the \texttt{AuditService} gracefully logs to `stderr` without crashing the core financial transaction. 

\subsection{Fault Tolerance \& Circuit Breakers}
Distributed systems face partial failures. ByteVault implements explicit timeouts when querying remote resources (like the Sub Branch via FDW).

\begin{lstlisting}[style=typescript]
// Snippet from backend/src/services/AccountService.ts
const foreign = await poolA().query(
  'SET statement_timeout = 2000; SELECT id FROM fdw_sub.accounts WHERE account_number = $1',
  [accountNumber]
);
\end{lstlisting}

\textbf{Why it is there:} If the \texttt{branch\_b\_db} container crashes, the FDW connection would normally hang indefinitely, consuming connection pool slots and eventually bringing down the Main API. The 2-second \texttt{statement\_timeout} acts as a Circuit Breaker, failing fast and keeping the system responsive.

\subsection{Centralized Error Handling}
To standardize API responses and prevent raw stack traces from leaking to the client, ByteVault uses a custom \texttt{AppError} utility.

\begin{lstlisting}[style=typescript]
// Snippet from backend/src/utils/errors.ts
export class AppError extends Error {
  constructor(public message: string, public status: number = 400) {
    super(message);
  }
  static notFound(msg: string) { return new AppError(msg, 404); }
}
\end{lstlisting}

The global Express error handler catches these and shapes them into predictable JSON structures: \texttt{\{ "error": "Message" \}}.

\subsection{Idempotency Keys}
All state-mutating payment and ledger routes require an \texttt{Idempotency-Key} HTTP header. 

\textbf{Why it is there:} If a client's network drops \textit{after} the server processes a transfer, the client might retry. Without idempotency, this results in a double-charge.

The middleware calculates a hash of the request payload and checks the \texttt{idempotency\_keys} table. 
\begin{itemize}
    \item \textbf{New Request}: Marks status \texttt{IN\_PROGRESS}.
    \item \textbf{Concurrent Request}: Blocks with 409 Conflict.
    \item \textbf{Retry}: Returns the cached \texttt{response\_body} immediately.
\end{itemize}

\subsection{Security \& Authentication}
ByteVault uses Stateless JWT (JSON Web Tokens) for employee authentication. 

\begin{lstlisting}[style=typescript]
export type JwtEmployeeClaims = {
  sub: string;
  email: string;
  role: EmployeeRole;
  branchId: string | null;
  branchName: string | null;
};
\end{lstlisting}

\textbf{Why it is there:} Stateless JWTs eliminate the need for server-side session stores, improving scalability. The JWT explicitly embeds the \texttt{branchName} and \texttt{role} (e.g., MAKER, ADMIN) so the frontend can quickly render UI elements without additional API calls.
